% ============================================================
\section{模型假设}

为简化问题、聚焦核心矛盾，本文做出以下假设：

\textbf{假设1：} 将进入日内调度与执行的负荷、光伏功率及电价序列统一离散到10 min时间尺度。在每个调度时段内，各变量采用该时段的代表值，忽略时段内部的短时波动；功率对应的时段电量按
\[
E_t=P_t\Delta t,\qquad \Delta t=\frac{1}{6}\text{ h}
\]
计算，购电费用按各类购电量与其对应结算单价计算。

\textbf{假设2：} 将小区微网等效为聚合能源系统，仅在系统总量层面考虑负荷、光伏发电、外网购电、储能充放电及弃光等能量流动；对于不同问题中涉及的调整购电、紧急购电等方式，均纳入相应时段的总能量平衡。忽略配电网内部线路损耗、节点电压及潮流分布等空间因素。

\textbf{假设3：} 对储能性能的时间变化作短期稳定化处理。在研究周期内，储能系统的容量边界、充放电功率上限及充放电效率按题目给定值和模型设定的固定值处理，不考虑温度变化、设备老化、自放电及故障等因素引起的变化。

\textbf{假设4：} 在每个调度或滚动决策时刻，模型规定可使用的历史观测、当前储电量及已发布预报能够及时、准确获得，且调度指令能够在相应时段内执行，忽略测量误差、通信延迟和控制执行滞后；对于决策时尚未知的未来信息，仍按照各问题所建立的预测或估计方法处理。
